\documentclass[11pt]{article}
\usepackage{amsmath,amssymb,amsthm,hyperref}
\newtheorem{proposition}{Proposition}
\title{An exact obstruction to extracting the triple term linearly}
\author{Exact certificate; independently verified by a second agent}
\date{30 September 2026}
\begin{document}
\maketitle

For a concatenation of $m$ independently specified palindrome blocks
$\rho_i\operatorname{rev}(\rho_i)$ on $n$ labels, write $R(\sigma)$ for
the ratio of the probability of target order $\sigma$ to $1/n!$.
Let $S(\sigma)$ be the triple contribution defined in the companion
palindrome notes. Both $R-1$ and $S$ are vectors indexed by $S_n$.

\begin{proposition}
Already for $n=5$ and $m=3$, there is no linear map $T$ on target-order
vectors such that $S=T(R-1)$ for every choice of the block permutations.
This remains false if all lower-dimensional ordering marginals are
included in the input, since they are linear averages of the full vector.
\end{proposition}

The conclusion concerns an exact universal linear extraction. It does
not exclude an approximate extraction, a nonlinear inequality, or a
lower bound for the complete deterministic ordering error. In this
example products of two block defects have degree at least six and
therefore do not occur: the obstruction is cancellation between
single-block contributions of lengths three, four, and five.

\begin{proof}[Proof by an exact integer certificate]
For a distinct-letter pattern $p$ of length $k\in\{3,4,5\}$ put
\[
 D_\rho(p)=N_\rho(p)-\frac{2^k}{k!},
\]
where $N_\rho(p)$ is its subsequence count in the ten-letter palindrome
$\rho\operatorname{rev}(\rho)$. Let slots be $i=0,1,2$ and write
$j=2-i$. Define integer columns, indexed by full patterns $p\in S_5$,
\[
 A_{i,\rho}(p)=15\sum_{\substack{a,b\ge0\\a+b\le2}}
 \frac{(2i)^a(2j)^b}{a!b!}\,
 D_\rho(p_{a+1}\cdots p_{5-b}),
\]
\[
 B_{i,\rho}(p)=15\sum_{\substack{a,b\ge0\\a+b=2}}
 \frac{(2i)^a(2j)^b}{a!b!}\,
 D_\rho(p_{a+1}\cdots p_{5-b}).
\]
The reduced-signature product expansion gives exactly
\[
 R-1=\frac{120}{15\cdot6^5}\sum_{i=0}^2 A_{i,\rho_i},\qquad
 S=\frac{120}{15\cdot6^5}\sum_{i=0}^2 B_{i,\rho_i}.
\]
Indeed, a palindrome block agrees with $\exp(2\sum X_i)$ through
degree two; products of two defects vanish in degree five. The formula
for $B$ retains precisely its degree-three part.

Here is a signed integer array $v_{i,\rho}$, with unlisted entries zero:
\[
\begin{array}{c|c|r@{\qquad}c|c|r}
 i&\rho&v&i&\rho&v\\\hline
 0&01234&-5&0&01324&5\\
 0&10234&-4&0&10324&4\\
 0&23401&8&0&32401&-8\\
 0&40213&-8&0&40312&8\\
 0&41203&-8&0&41302&8\\
 1&01234&2&1&01324&-2\\
 2&01234&-1&2&01324&1
\end{array}
\]
Exact subsequence counting verifies
\[
 \sum_\rho v_{i,\rho}=0\quad(i=0,1,2),\qquad
 \sum_{i,\rho}v_{i,\rho}A_{i,\rho}=0,
\]
whereas
\[
 \left(\sum_{i,\rho}v_{i,\rho}B_{i,\rho}\right)(01234)=-120.
\]
All 120 coordinates of the first vector identity and of the nonzero
second vector are recorded in
\path{linear_extraction_counterexample.json}. Two independent exact
programs verified them: one uses subsequence dynamic programming, the
other enumerates position subsets of each palindrome directly.

To deduce the stated impossibility for actual block choices, independently
sample each slot from the rational probability distribution
\[
 q_{i,\rho}=\frac1{120}+\frac{v_{i,\rho}}{960}.
\]
These probabilities are nonnegative because $|v|\le8$, and each slot
sums to one. Under the unperturbed uniform distributions, both displayed
defect vectors have mean zero, by uniform relabeling. The exact identities
therefore give
\[
 \mathbb E(R-1)=0,\qquad
 \mathbb E S(01234)=-\frac1{6^5}\ne0.
\]
A universal linear map $T$ would imply
$\mathbb ES=T(\mathbb E(R-1))=0$, a contradiction.
\end{proof}

\paragraph{Reproduction.}
The discovery script is \path{find_extraction_counterexample.py}; the
matrix definition is \path{check_linear_extraction.py}. The independent
standard-library verifier is
\path{../independent_directions/review_linear_extraction_counterexample.py}.
The exact identities, not the numerical ranks used during discovery,
are the certificate.
\end{document}
